\documentclass[11pt]{article}

\usepackage[margin=1in]{geometry}
\usepackage{booktabs}
\usepackage{amsmath}
\usepackage{amssymb}
\usepackage{hyperref}
\usepackage{cite}
\usepackage{authblk}

\title{An Atlas of Brain--Bone Sympathetic Neural Circuits}

\author[1]{Anonymous Author}
\affil[1]{Anonymous Institution}

\date{}

\begin{document}

\maketitle

\begin{abstract}
There is clear evidence that the sympathetic nervous system (SNS) mediates bone metabolism: histological studies show abundant SNS innervation of the periosteum and bone marrow with noradrenergic fibers that immunostain for tyrosine hydroxylase, dopamine beta hydroxylase, or neuropeptide Y. Nonetheless, the brain sites that send efferent SNS outflow to bone have not yet been characterized. Using pseudorabies virus (PRV152) transneuronal tracing from the right murine femur ($N=6$ male mice, 3--4 months old; microinjections $4.7 \times 10^9$\,pfu/mL at five loci, $150$\,nL/locus), we report, for the first time, the identification of central SNS outflow sites that innervate bone. We identify $87$ PRV152-positive brain nuclei, sub-nuclei and regions distributed across six brain divisions. Neuron counts (mean $\pm$ SEM over the four mice with comparable infection) were greatest in hypothalamus ($1177.25 \pm 62.75$), followed by midbrain and pons ($1065 \pm 22.39$), hindbrain medulla ($495.25 \pm 33.49$), forebrain ($237.5 \pm 15.08$), cerebral cortex ($104.75 \pm 4.64$), and thalamus ($65.25 \pm 7.78$). Certain sites, including the raphe magnus (RMg) of the medulla and the periaqueductal gray (PAG) of the midbrain, display disproportionately high PRV152 infection, suggesting site-specific variation in central SNS outflow to bone. This atlas of brain nodes controlling SNS efferent signals to the femur complements prior atlases of SNS outflow to bone marrow and adipose depots, and it provides an anatomical scaffold for studying the neural regulation of bone metabolism and the mechanisms of central bone pain.
\end{abstract}

\section{Introduction}

Elegant studies have established that increased SNS tone causes bone loss through reduced bone formation coupled with increased bone resorption \cite{elefteriou2018,elefteriou2005,takeda2002}. Leptin-deficient mice display a high bone-mass phenotype, and the anti-osteogenic action of leptin is mediated centrally by glucose-responsive neurons in the ventromedial hypothalamus that recruit peripheral SNS pathways \cite{ducy2000,takeda2002}. Both the periosteum and the bone marrow are richly innervated by SNS fibers immunoreactive for tyrosine hydroxylase, dopamine beta hydroxylase, or neuropeptide Y, largely associated with vasculature and with SNS vesicular acetylcholine transporter (VAChT), whereas vasoactive intestinal polypeptide (VIP) fibers are mainly parenchymal \cite{francis1997,hill1991,hohmann1986,martin2007}. Despite this, the distribution of SNS nerves within the mammalian skeleton and their connectivity to central neurons remains incompletely understood.

Viral transneuronal tracing is the established technology for defining central SNS circuitry to peripheral organs. Bartha's K-strain of pseudorabies virus (PRV) is a transneuronal tract tracer that maps multi-synaptic circuits within one animal \cite{ekstrand2008,enquist2002,song2005a}. PRV is endocytosed at axon terminals, transported retrogradely, and then jumps across synapses to higher CNS sites \cite{curanovic2009}. Using this technology, previous work defined postganglionic SNS innervation of white and brown adipose tissue depots with separate and shared central SNS relay sites \cite{ryu2014,ryu2015}, and a direct neuroanatomical connection between phosphodiesterase 5A (PDE5A)-containing neurons in specific brain nuclei and bone \cite{kim2020pde5a}. A hierarchical circuit controlling SNS output to rat femoral epiphyseal bone marrow has also been defined \cite{denes2005}, identifying PRV-labelled neurons in C1, A5, A7 catecholaminergic groups, the ventrolateral and ventromedial medulla, periaqueductal gray, the paraventricular hypothalamic nucleus and other hypothalamic nuclei, as well as the insular and piriform cortex. However, no study has mapped the localization and organization of the central SNS circuitry innervating the \emph{murine femur}. Here we address this gap.

\section{Methods}

\paragraph{Animals.} Adult male mice ($\sim$3 to 4 months old) were single-housed under a 12\,h:12\,h light:dark cycle at $22 \pm 2\,^\circ\mathrm{C}$ with ad libitum access to water and regular chow. All procedures were approved by the Mount Sinai Institutional Animal Care and Use Committee and performed in accordance with PHS and USDA guidelines.

\paragraph{Viral tracing.} To identify brain sites sending SNS outflow to bone, we used the transsynaptic tracer PRV152, which expresses enhanced green fluorescent protein (EGFP) under control of the human cytomegalovirus immediate-early promoter. All virus injections followed Biosafety Level 2 standards. Mice ($N = 6$) were anesthetized with isoflurane (2--3\% in oxygen; Baxter Healthcare, Deerfield, IL) and the right femur--tibia joint was exposed for a series of PRV152 microinjections ($4.7 \times 10^9$\,pfu/mL) at five loci ($150$\,nL/locus) evenly distributed across the bone metaphysis and periosteum, which are known to be enriched with SNS innervation. The syringe was held in place for $60$\,s after each injection to prevent efflux of virus. The incision was closed with sterile sutures and wound clips. Nitrofurazone powder (nfz Puffer; Hess \& Clark, Lexington, KY) was applied locally to minimize the risk of bacterial infection. As a control, PRV152 placed on the bone surface yielded no EGFP signal in the paraventricular hypothalamic nucleus (PVH) or raphe pallidus (RPa).

\paragraph{Histology.} Animals were sacrificed $6$ days after the last PRV152 injection, based on the progression of the virus to the brain in pilot studies. Mice were euthanized with carbon dioxide and perfused transcardially with $0.9\%$ heparinized saline followed by $4\%$ paraformaldehyde in $0.1$\,M phosphate-buffered saline (PBS; pH $7.4$). Brains were post-fixed in the same fixative for $3$ to $4$ hours at $4\,^\circ\mathrm{C}$, then transferred to a $30\%$ sucrose solution in $0.1$\,M PBS with $0.1\%$ sodium azide and stored at $4\,^\circ\mathrm{C}$ until sectioning at $25\,\mu\mathrm{m}$ on a freezing-stage sliding microtome. Sections were stored in $0.1$\,M PBS with $0.1\%$ sodium azide until immunofluorescence.

\paragraph{Immunofluorescence.} Free-floating brain sections were rinsed in $0.1$\,M PBS ($2 \times 15$\,min), blocked for $30$\,min in $10\%$ normal goat serum (NGS; Vector Laboratories, Burlingame, CA) and $0.4\%$ Triton X-100 in $0.1$\,M PBS, and then incubated in primary chicken anti-EGFP ($1{:}1000$; Thermo Fisher Scientific, catalog no.\ A10262) for $18$\,hours. Sections were subsequently incubated in AlexaFluor-488-coupled goat anti-chicken secondary ($1{:}700$; Jackson Immunoresearch, catalog no.\ 103-545-155) with $2\%$ NGS and $0.4\%$ Triton X-100 in $0.1$\,M PBS at room temperature for $2$\,hours. Controls either omitted the primary antibody or pre-adsorbed it with the immunizing peptide overnight at $4\,^\circ\mathrm{C}$, yielding no immunoreactive staining. Sections were mounted with ProLong Gold Antifade (Thermo Fisher Scientific, catalog no.\ P36982). Images were acquired at $10\times$ and $20\times$ using an Observer.Z1 fluorescence microscope (Carl Zeiss) and overlaid with DAPI in Zen (Zeiss) and ImageJ (NIH). Scale bars $= 50\,\mu\mathrm{m}$.

\paragraph{Quantification.} PRV152-immunoreactive cells were counted manually in every sixth brain section using the tag feature of Adobe Photoshop CS5.1 to avoid double counting. Neuron numbers were averaged across each nucleus, sub-nucleus and region per animal, and nuclei were identified against the Paxinos and Franklin mouse brain atlas \cite{paxinos2007}.

\section{Results}

\subsection{Validation of PRV152 retrograde infection}
Mice remained asymptomatic until day $5$ post-inoculation, after which they displayed occasional body-weight loss, decreased mobility, and an ungroomed coat. Four of the six PRV152-injected mice exhibited equivalent infection across the hindbrain-to-forebrain neuroaxis and were included in analyses; two mice showed over-infection (cloudy plaques surrounding infected neurons) and were excluded. Unilateral PRV152 injection into the right femur produced bilateral brain labelling without noticeable hemispheric dominance. Control PRV152 application onto the bone surface, rather than injection into periosteum or metaphysis, produced no EGFP signal in PVH or RPa, whereas periosteum or metaphysis injections produced positive EGFP in PVH, and PRV152-labelled neurons were also detected in the intermediolateral cell column (IML) of the spinal cord at T13--L2 levels, consistent with a bone--SNS ganglia--IML--brain route of infection.

\subsection{Distribution of PRV152 labelling across six brain divisions}
We identified $87$ PRV152-positive brain nuclei, sub-nuclei and regions across six brain divisions. Neuron counts per division (mean $\pm$ SEM) are summarized in Table~\ref{tab:divisions}: hypothalamus contained the most infected SNS neurons ($1177.25 \pm 62.75$), followed by midbrain and pons ($1065 \pm 22.39$), hindbrain medulla ($495.25 \pm 33.49$), forebrain ($237.5 \pm 15.08$), cerebral cortex ($104.75 \pm 4.64$), and thalamus ($65.25 \pm 7.78$).

\begin{table}[h]
\centering
\caption{PRV152-labelled neuron counts (mean $\pm$ SEM) across the six brain divisions that send SNS outflow to the femur.}
\label{tab:divisions}
\begin{tabular}{lc}
\toprule
Brain division & PRV152-labelled neurons \\
\midrule
Hypothalamus & $1177.25 \pm 62.75$ \\
Midbrain and pons & $1065 \pm 22.39$ \\
Hindbrain medulla & $495.25 \pm 33.49$ \\
Forebrain & $237.5 \pm 15.08$ \\
Cerebral cortex & $104.75 \pm 4.64$ \\
Thalamus & $65.25 \pm 7.78$ \\
\bottomrule
\end{tabular}
\end{table}

\subsection{Dominant nuclei within each division}
Within the hypothalamus, the lateral hypothalamus (LH), paraventricular hypothalamic nucleus (PVH), and dorsomedial hypothalamus (DM) showed the highest percentages of PRV152-labelled neurons, with the LH and PVH also among the regions with the highest absolute numbers of the $25$ PRV152-positive hypothalamic nuclei, sub-nuclei and regions. In the midbrain and pons, the periaqueductal gray (PAG), lateral PAG (LPAG), and pontine reticular nucleus oral part (PnO) showed the highest PRV152 labelling across $18$ nuclei, sub-nuclei and regions. In the hindbrain medulla, the raphe pallidus (RPa), raphe magnus (RMg), and gigantocellular reticular nucleus (Gi) were among $23$ PRV152-positive nuclei, sub-nuclei and regions, and carried the largest percentages and counts of PRV152-labelled neurons. In the forebrain, the medial preoptic nucleus medial part (MPOM), bed nucleus of the stria terminalis (BST), and lateral septal nucleus ventral part (LSV) were the dominant nuclei among $15$ PRV152-positive forebrain sites. The cerebral cortex contained only three PRV152-labelled regions: the primary somatosensory cortex hindlimb region (S1HL) and the secondary and primary motor cortices (M2 and M1), with S1HL and M2 having the highest percentages and counts. The thalamus contained three PRV152-infected sites, with the periventricular fiber system (pv) and the precommissural nucleus (PrC) possessing the highest percentages and counts (Table~\ref{tab:topsites}).

\begin{table}[h]
\centering
\caption{Dominant nuclei, sub-nuclei and regions within each brain division (number of PRV152-positive sites per division shown in parentheses).}
\label{tab:topsites}
\begin{tabular}{ll}
\toprule
Division (PRV152-positive sites) & Top-ranked nuclei \\
\midrule
Hypothalamus (25) & LH, PVH, DM \\
Midbrain and pons (18) & PAG, LPAG, PnO \\
Hindbrain medulla (23) & RPa, RMg, Gi \\
Forebrain (15) & MPOM, BST, LSV \\
Cerebral cortex (3) & S1HL, M2, M1 \\
Thalamus (3) & pv, PrC \\
\bottomrule
\end{tabular}
\end{table}

\section{Discussion}

Using PRV152 retrograde transneuronal tracing from the murine femur, we report the first comprehensive atlas of central SNS outflow to bone. We identified $87$ brain nuclei, sub-nuclei and regions across six brain divisions that send SNS efferents to the femur. The atlas confirms and substantially extends prior work on SNS outflow to rat femoral epiphyseal bone marrow \cite{denes2005}, PDE5A-containing neurons projecting to bone \cite{kim2020pde5a}, and central SNS circuits to adipose depots \cite{ryu2014,ryu2015,ryu2017,bamshad1998,bamshad1999}. It demonstrates that bone is not innervated by a unique neuron group but rather by overlapping SNS circuitry common to the control of other peripheral targets.

Several nodes are particularly prominent. The hypothalamus contains the largest number of PRV152-labelled neurons ($1177.25 \pm 62.75$), with LH, PVH and DM dominating. The PVH is a classical source of SNS pre-autonomic neurons, whereas the LH predominantly innervates the femur in our dataset. Given that LH, PVH and DM all express leptin receptors \cite{flak2016leptin}, the anti-osteogenic relay for leptin \cite{ducy2000,takeda2002} may originate from neurons in these nuclei. The midbrain and pons division ($1065 \pm 22.39$) is dominated by the PAG, which receives afferent fibres from the parabrachial nucleus and RPa \cite{mantyh1982} and from the spinal cord \cite{pechura1986}, and which is responsible for descending modulation of pain perception \cite{baptista2018,benarroch2008,calvino2006}. The PAG is therefore well positioned to integrate bone-sensory inflow with SNS outflow. The medulla ($495.25 \pm 33.49$) is dominated by RPa and RMg: RPa projections to dorsal spinal horn regulate enkephalin release and opiate-mediated antinociception \cite{francois2017}, while RMg neurons modulate noxious stimuli \cite{fields1991}. Together the RMg--PAG axis may form part of an ascending hierarchical circuit relating to perception of bone pain.

Comparisons with prior tracing from rat femoral bone marrow \cite{denes2005} reveal extensive overlap in central nodes (PAG, somatosensory cortex, MPOM, thalamic periventricular nucleus, PVH, LH, RPa), suggesting that largely the same brain sites project to both femur and femoral bone marrow. However, the different relative infection levels across nuclei between the two targets argue for separable, site-specific SNS circuits within this shared architecture. The identification of dopamine-transporter-deficient osteopenia \cite{bliziotes2000}, the anti-osteogenic action of leptin \cite{takeda2002,ducy2000}, and the SNS innervation of bone \cite{francis1997,hill1991,hohmann1986,martin2007} are collectively consistent with our atlas.

Limitations of this study include the small cohort of four analyzable mice, the use of male mice only, and the single-target (femur) design. Future work should extend this atlas across sex, age and other skeletal sites, and should combine it with functional manipulations (e.g.\ chemogenetic silencing of LH, PVH, PAG, or RMg/RPa) to test causal roles in bone remodelling and bone pain. The overlap between our atlas and the SNS relay to adipose depots \cite{ryu2014,ryu2015,bamshad1998,bamshad1999,bowers2004,shi2001,song2001,song2005b,song2008,vaughan2012,leitner2009,nguyen2014,adler2012} suggests that future work could test whether manipulation of shared nodes differentially affects bone, bone marrow, and adipose tissue.

In summary, this atlas of brain--bone SNS circuitry provides an anatomical foundation for investigating the neural regulation of bone metabolism and of central bone pain, and it pinpoints actionable targets --- notably the RMg--PAG axis, PVH, LH, and S1HL --- for future mechanistic and therapeutic studies.

\bibliographystyle{unsrt}
\bibliography{references}

\end{document}
